%% =====================================================================
%%  main.tex
%%  Archivo maestro del informe. Compile con:  make
%%  (o bien: latexmk -xelatex main.tex)
%% =====================================================================
\documentclass[
  fuente=auto,              % auto | arial | liberation
  alineacion=justificada,   % izquierda (APA) | justificada
  paginacion=institucional, % institucional | apa
  interlineado=2,           % 1.5 | 2
  capitulos=arabigos,       % romanos | arabigos | sin
  portada=institucional     % institucional | apa | ninguna
]{si3apa7}

\usepackage{entornos-apa}

%% =====================================================================
%%  metadatos.tex
%%  Datos del informe. Edite únicamente este archivo para actualizar
%%  la portada, los encabezados y los metadatos del PDF.
%% =====================================================================

%% --- Datos institucionales -------------------------------------------
\newcommand{\universidad}{Universidad Privada Domingo Savio}
\newcommand{\sede}{Sede Potosí}
\newcommand{\facultad}{Facultad de Ingeniería}
\newcommand{\carrera}{Carrera de Ingeniería de Sistemas}
\newcommand{\filiacion}{Universidad Privada Domingo Savio, Facultad de Ingeniería}
\newcommand{\ciudad}{Potosí, Bolivia}

%% --- Datos del documento ---------------------------------------------
\newcommand{\tipotrabajo}{Informe de Proyecto de Grado}
\newcommand{\materia}{Sistemas de Información III}
\newcommand{\titulo}{Sistema web para la gestión de paquetes turísticos y reservas en Vizcacha Tours, Potosí}
\newcommand{\subtitulo}{Análisis y propuesta de un sistema de información transaccional}

%% --- Datos de los estudiantes y del docente ------------------------------
%%  El trabajo tiene dos autores. Cada uno se declara por separado;
%%  \autor resume la lista para los metadatos del PDF.
\newcommand{\autorUno}{Roberto Carlos Ereño Acuña}
\newcommand{\autorDos}{Erick Fernando Velasco Flores}
\newcommand{\autor}{\autorUno\space y\space \autorDos}
\newcommand{\correo}{estudiante@upds.edu.bo}
\newcommand{\docente}{Ing. Carmen Nelinda Ozuna Villca}

%% --- Datos de entrega ------------------------------------------------
\newcommand{\gestion}{2026}
\newcommand{\mesentrega}{Noviembre}
\newcommand{\fechaentrega}{16 de noviembre de 2026}

%% --- Bibliografía ----------------------------------------------------
\addbibresource{informe/bibliografia/referencias.bib}

%% =====================================================================
%%  portada.tex
%%  Portada institucional UPDS y portada APA 7 alternativa.
%%  Los datos se editan en config/metadatos.tex.
%% =====================================================================

%% ---------------------------------------------------------------------
%%  Portada institucional (predeterminada)
%% ---------------------------------------------------------------------
\newcommand{\portadaInstitucional}{%
  \begin{titlepage}%
  \estiloPortada
  \singlespacing
  \begin{center}
    \IfFileExists{informe/figuras/logos/upds.png}%
      {\includegraphics[width=4.6cm]{informe/figuras/logos/upds.png}\par}%
      {\fbox{\parbox[c][2cm][c]{5cm}{\centering \small Logotipo UPDS\\ (informe/figuras/logos/upds.png)}}\par}
    \vspace{0.9cm}
    {\bfseries\large \universidad\par}
    \vspace{0.2cm}
    {\bfseries\normalsize \facultad\par}
    {\bfseries\normalsize \carrera\par}
    {\normalsize \sede\par}
    \vspace{1.3cm}
    {\bfseries\large \tipotrabajo\par}
    \vspace{0.2cm}
    {\normalsize Materia: \materia\par}
    \vspace{1.1cm}
    {\bfseries\Large \titulo\par}
    \vspace{0.4cm}
    {\normalsize\itshape \subtitulo\par}
    \vspace{1.2cm}
    {\normalsize Presentado por:\par}
    \vspace{0.15cm}
    {\bfseries\normalsize \autor\par}
    {\normalsize Registro universitario: \registro\par}
    \vspace{1.0cm}
    {\normalsize Docente:\par}
    \vspace{0.15cm}
    {\bfseries\normalsize \docente\par}
    \vfill
    {\normalsize \ciudad\par}
    {\bfseries\normalsize \gestion\par}
  \end{center}
  \end{titlepage}%
}

%% ---------------------------------------------------------------------
%%  Portada APA 7 de estudiante (opcional: portada=apa)
%% ---------------------------------------------------------------------
\newcommand{\portadaAPA}{%
  \begin{titlepage}%
  \estiloPortada
  \singlespacing
  \begin{center}
    \vspace*{3.5cm}
    {\bfseries\large \titulo\par}
    \vspace{2.5cm}
    {\normalsize \autor\par}
    {\normalsize \filiacion\par}
    {\normalsize \materia\par}
    {\normalsize Docente: \docente\par}
    {\normalsize \fechaentrega\par}
  \end{center}
  \end{titlepage}%
}


\begin{document}

%% ------------------------- Preliminares ------------------------------
\portadaDocumento
\iniciarPreliminares

\preliminar{Resumen}

El presente proyecto aborda la ineficiencia del proceso de venta y reserva
de paquetes turísticos en Vizcacha Tours, Potosí, donde la verificación de
cupos y la actualización de la información no se realizan de manera uniforme
entre el personal. El objetivo general es desarrollar un sistema web de
gestión de reservas que mejore la consistencia del proceso de venta y reserva
de paquetes turísticos en Vizcacha Tours. La propuesta se sustenta en los
antecedentes de los sistemas de reservas, en el marco de la gestión de
paquetes turísticos y en criterios de usabilidad y seguridad para
aplicaciones web. La metodología prevé entrevistas, observación y revisión
autorizada de documentos en la agencia, con el fin de diagnosticar el proceso
actual y diseñar la solución de acuerdo con sus necesidades. El sistema será
utilizado por el personal de la agencia y por los clientes, que accederán con
una cuenta propia para consultar la disponibilidad, solicitar reservas y
registrar sus pagos; ambas entradas siguen el mismo criterio de validación y
de confirmación. Se espera que el sistema contribuya a reducir inconsistencias
en la disponibilidad, disminuir el tiempo de confirmación y ofrecer
información confiable para la toma de decisiones.

\vspace{0.5\baselineskip}
\noindent\textbf{Palabras clave:} agencias de viaje, paquetes turísticos,
reservas, sistemas web, control de disponibilidad


\clearpage
\tableofcontents

\clearpage
\phantomsection
\addcontentsline{toc}{chapter}{Índice de tablas}
\listoftables

\clearpage
\phantomsection
\addcontentsline{toc}{chapter}{Índice de figuras}
\listoffigures

%% ---------------------------- Cuerpo ---------------------------------
\iniciarCuerpo

\capituloSinNumero{Introducción}
\label{cap:introduccion}

El turismo es una actividad económica relevante para Potosí. La ciudad
histórica fue inscrita en la Lista del Patrimonio Mundial en 1987
\parencite{unescoPotosi}, y el Instituto Nacional de Estadística publica
información sobre el flujo de visitantes internacionales en Bolivia
\parencite{ineVisitantes}. En este contexto, las agencias de turismo de la
ciudad ofrecen paquetes que combinan transporte, alojamiento, guías y
actividades; la información sobre precios, itinerarios, fechas de salida y
cupos debe estar disponible para el personal y para los clientes.

La venta de estos servicios evolucionó desde la atención presencial y
telefónica, apoyada en catálogos impresos, hacia canales digitales como
sitios web, redes sociales, mensajería instantánea y plataformas de reserva.
Los sistemas globales de distribución (GDS) permitieron consultar horarios,
tarifas y disponibilidad de distintos proveedores desde un mismo punto de
acceso \parencite{amadeusGDS}; más tarde, plataformas como Booking.com
ampliaron la consulta y la reserva en línea \parencite{booking}. En Potosí,
las agencias suelen combinar la atención en oficina con estos canales, por lo
que la información debe mantenerse coherente entre todos ellos.

En el mercado existen soluciones consolidadas para el sector turístico.
Amadeus ofrece servicios de distribución y reserva conectados con múltiples
proveedores; Booking.com facilita la publicación de alojamientos y la reserva
en línea; y Ezus reúne cotizaciones, itinerarios, tarifas y documentación para
agencias de viajes \parencite{ezus}. Estas herramientas aportan
centralización de la información, consulta de disponibilidad, trazabilidad de
las operaciones y atención por varios canales. En el ámbito regional,
\textcite{huaman2017} desarrolló un sistema de reservas para la agencia Rap
Travel.

El proyecto se orienta a Vizcacha Tours, una agencia de Potosí que administra
paquetes propios. A diferencia de las plataformas internacionales, no
pretende conectarse automáticamente con aerolíneas ni con pasarelas
bancarias, sino organizar la información de los paquetes, las salidas y las
reservas, y mantener un control consistente de los cupos. La comparación con
las soluciones existentes permite identificar funciones útiles y definir un
alcance adecuado al tamaño y a las necesidades de la agencia.

\capituloSinNumero{Antecedentes}
\label{cap:antecedentes}

Antes de la expansión de Internet, muchas agencias de turismo atendían las
consultas directamente en sus oficinas o por teléfono. Los catálogos
impresos y los contactos con empresas de transporte y hospedaje servían para
presentar las ofertas. Con estos medios era posible trabajar, siempre que se
mantuviera un control adecuado de los cambios y las confirmaciones.

Posteriormente se extendió el uso de sistemas computarizados de reservas. Los
sistemas globales de distribución, conocidos como GDS, facilitaron la
consulta de horarios, tarifas y disponibilidad de diferentes proveedores
turísticos desde un punto de acceso común \parencite{amadeusGDS}. Esta
evolución permitió organizar mejor la información de servicios que antes se
consultaban por separado.

Con el crecimiento de los servicios en línea, los clientes comenzaron a
buscar y reservar opciones de viaje desde sus computadoras y celulares.
Booking.com, cuya actividad comenzó en 1996, es un ejemplo de la expansión de
las plataformas digitales en el sector turístico \parencite{booking}. Una
agencia de Potosí que ofrece paquetes propios no necesita copiar todas las
funciones de una plataforma internacional; lo principal es que pueda
administrar correctamente sus servicios y cupos.

En los últimos años se han desarrollado herramientas para elaborar
cotizaciones, organizar itinerarios, guardar información de proveedores y
controlar pagos. Por ejemplo, Ezus presenta una solución que reúne
cotizaciones, itinerarios, tarifas y documentación \parencite{ezus}. Estas
funciones muestran cómo distintos procesos se pueden relacionar dentro de un
mismo sistema, aunque cada agencia debe seleccionar lo que realmente
necesita.

\section*{Relación con sistemas similares}

Como referencia internacional se puede mencionar a Amadeus, que facilita el
acceso a información de servicios turísticos de distintos proveedores, y a
Booking.com, que permite consultar y reservar opciones de viaje en línea. Por
otra parte, programas de gestión como Ezus incluyen herramientas
administrativas para preparar propuestas, organizar itinerarios y manejar la
información comercial de una agencia \parencite{amadeusGDS,booking,ezus}.

En la región existe un trabajo cercano a este
proyecto. \textcite{huaman2017} desarrolló un sistema de información para la
gestión de reservas de paquetes turísticos en la agencia Rap Travel, tras
identificar que el proceso se realizaba manualmente y sin un control adecuado
del seguimiento de los pagos. En la búsqueda documental realizada para el
proyecto no se encontró documentación pública sobre sistemas de gestión de
reservas para agencias de turismo de Potosí, por lo que el diagnóstico en
Vizcacha Tours permitirá documentar su situación.

El proyecto comparte con estos sistemas la idea de mantener la información
organizada y actualizada, pero tendrá un alcance menor. No se pretende crear
una plataforma mundial ni conectarse directamente con aerolíneas: el sistema
se enfocará en los paquetes que ofrece Vizcacha Tours y en la forma de
consultar, reservar y controlar los cupos de cada salida programada. También
se tomarán en cuenta criterios básicos de usabilidad, pues la norma ISO
9241-210:2019 destaca la importancia de diseñar sistemas interactivos
considerando las necesidades de las personas usuarias \parencite{iso9241}.

\section*{Flujo de información de la agencia}

Vizcacha Tours ofrece información sobre destinos, paquetes, itinerarios,
precios, fechas, cupos, servicios incluidos y condiciones de venta. Esa
información se difunde en la oficina, por teléfono, mediante mensajería
instantánea, en redes sociales y, cuando corresponde, a través de plataformas
digitales de reserva. Los clientes pueden consultar, comparar y solicitar una
propuesta por cualquiera de esos canales. Con el sistema propuesto, el acceso
web del cliente se convierte en un canal adicional de consulta, de modo que la
información de disponibilidad y de las reservas se mantiene al día en un mismo
punto de acceso para todos los canales de atención.

La información que circula incluye consultas, datos de clientes y pasajeros,
cotizaciones, reservas, pagos y comprobantes, además de la coordinación con
proveedores de transporte, hospedaje y guías. El flujo general comienza con la
recepción de la consulta, continúa con el registro y la propuesta, y exige que
las áreas mantengan actualizados los datos de disponibilidad y las
condiciones del servicio. La consulta y la solicitud de reserva pueden
originarse en la atención del personal o directamente en el sistema, pero en
ambos casos deben seguir el mismo criterio de disponibilidad y de registro, y
las solicitudes que no se confirman en el plazo previsto se liberan de manera
automática. La descripción detallada del proceso de negocio se desarrolla más
adelante; en esta sección solo se presenta el flujo general de información y
los canales de atención.

\capituloSinNumero{Situación problemática}
\label{cap:situacion}

La situación problemática se concentra en la ineficiencia del proceso de
venta de paquetes turísticos en Vizcacha Tours, Potosí, que afecta
principalmente el control de los cupos disponibles y la confirmación de las
reservas. El árbol de problemas que sintetiza las causas y los efectos de
esta situación se presenta en la figura~\ref{fig:arbol}.

Entre las posibles causas están la falta de un criterio uniforme para revisar
la disponibilidad antes de confirmar, la ausencia de reglas claras para
modificar los estados de reserva y la demora en comunicar cambios de
capacidad o fechas. Estos factores no dependen necesariamente del uso de
registros manuales o digitales: pueden presentarse si el proceso no está bien
definido.

De confirmarse estas dificultades, podrían ocasionar reservas mayores a la
capacidad de una salida, reprogramaciones, tiempo adicional de revisión y
reclamos por parte de los turistas. Además, la información de cotizaciones,
pagos o proveedores podría quedar descoordinada respecto de las reservas
confirmadas.

\begin{landscape}
\thispagestyle{paisaje}
\begin{figuraapa}{fig:arbol}{Árbol de problemas del proceso de venta de paquetes turísticos}
\centering
\scalebox{0.85}{%% =====================================================================
%%  arbol-problemas.tex
%%  Diagrama de árbol de problemas (TikZ).
%%  Se incluye desde el Anexo A en una página horizontal.
%% =====================================================================
\begin{tikzpicture}[
  nodo/.style={
    draw=black!65, line width=0.5pt, rounded corners=1pt,
    align=center, font=\scriptsize, text width=4cm,
    inner sep=3pt, minimum height=1.7cm, fill=white
  },
  central/.style={
    nodo, text width=17cm, font=\small\bfseries,
    fill=black!8, minimum height=1.3cm
  },
  flecha/.style={
    -{Latex[length=1.8mm,width=1.4mm]},
    draw=black!65, line width=0.5pt
  },
  rotulo/.style={
    rotate=90, font=\tiny\bfseries, text=black!70, align=center
  }
]

%% ------------------------- Efectos indirectos ------------------------
\node[nodo] (I1) at (-8.4,0) {Pérdida potencial de clientes y oportunidades de venta};
\node[nodo] (I2) at (-2.8,0) {Disminución de la confianza en el servicio ofrecido};
\node[nodo] (I3) at ( 2.8,0) {Mayores costos de coordinación y corrección};
\node[nodo] (I4) at ( 8.4,0) {Dificultades en la planificación de salidas turísticas};

%% -------------------------- Efectos directos -------------------------
\node[nodo] (D1) at (-8.4,-2.6) {Demoras en la confirmación de reservas};
\node[nodo] (D2) at (-2.8,-2.6) {Errores en las reservas, cupos duplicados o mal asignados};
\node[nodo] (D3) at ( 2.8,-2.6) {Información contradictoria sobre la disponibilidad};
\node[nodo] (D4) at ( 8.4,-2.6) {Reprogramaciones, reclamos y cancelaciones};

%% --------------------------- Problema central ------------------------
\node[central] (C) at (0,-5.2) {Ineficiencia en el proceso de venta de paquetes turísticos en Vizcacha Tours, Potosí};

%% --------------------------- Causas directas -------------------------
\node[nodo] (CD1) at (-8.4,-7.8) {Verificación no uniforme de cupos antes de confirmar una reserva};
\node[nodo] (CD2) at (-2.8,-7.8) {Desfase en la actualización de ventas entre distintos canales de atención};
\node[nodo] (CD3) at ( 2.8,-7.8) {Cambios de fechas o capacidad que no se reflejan oportunamente en las reservas};
\node[nodo] (CD4) at ( 8.4,-7.8) {Criterios poco claros para modificar, cancelar y liberar reservas};

%% ------------------------- Causas indirectas -------------------------
\node[nodo] (CI1) at (-8.4,-10.4) {Responsabilidades de verificación y actualización sin definición suficiente};
\node[nodo] (CI2) at (-2.8,-10.4) {Información de salidas y reservas sin una fuente compartida y actualizada};
\node[nodo] (CI3) at ( 2.8,-10.4) {Falta de controles para identificar cambios y prevenir conflictos de capacidad};
\node[nodo] (CI4) at ( 8.4,-10.4) {Procedimientos de seguimiento y validación insuficientemente estandarizados};

%% ------------------------------ Flechas ------------------------------
%% Efecto directo -> efecto indirecto
\foreach \d/\i in {D1/I1,D2/I2,D3/I3,D4/I4} {
  \draw[flecha] (\d.north) -- (\i.south);
}

%% Problema central -> efectos directos
\foreach \d/\x in {D1/-8.4,D2/-2.8,D3/2.8,D4/8.4} {
  \draw[flecha] (C.north -| \x,0) -- (\d.south);
}

%% Causas directas -> problema central
\foreach \c/\x in {CD1/-8.4,CD2/-2.8,CD3/2.8,CD4/8.4} {
  \draw[flecha] (\c.north) -- (C.south -| \x,0);
}

%% Causas indirectas -> causas directas
\foreach \c/\d in {CI1/CD1,CI2/CD2,CI3/CD3,CI4/CD4} {
  \draw[flecha] (\c.north) -- (\d.south);
}

%% --------------------------- Rótulos de nivel -------------------------
\node[rotulo] at (-10.95,0)     {E. INDIRECTOS};
\node[rotulo] at (-10.95,-2.6)  {E. DIRECTOS};
\node[rotulo] at (-10.95,-5.2)  {PROBLEMA};
\node[rotulo] at (-10.95,-7.8)  {C. DIRECTAS};
\node[rotulo] at (-10.95,-10.4) {C. INDIRECTAS};

\end{tikzpicture}
}
\vspace{2pt}
\begin{minipage}{0.88\linewidth}
  \nota{Fuente: elaboración propia a partir del diagnóstico preliminar.}
\end{minipage}
\end{figuraapa}
\end{landscape}

\capituloSinNumero{Problema científico}
\label{cap:problema-cientifico}

¿Cómo mejorar la consistencia del proceso de reservas de paquetes turísticos
en Vizcacha Tours de la ciudad de Potosí durante la gestión 2026?

\capituloSinNumero{Objeto de estudio y campo de acción}
\label{cap:objeto}

\section*{Objeto de estudio}

El objeto de estudio es el sistema de información transaccional (TPS) que
sustenta la operación comercial de Vizcacha Tours. Este tipo de sistema
registra y procesa las transacciones diarias de consulta, cotización, reserva
y pago, y organiza los datos de paquetes, salidas, clientes y pasajeros para
que el personal disponga de información actualizada. El sistema atiende a dos
grupos de usuarios: el personal de la agencia, que administra la información y
confirma las reservas, y los clientes, que acceden con una cuenta propia para
consultar la disponibilidad, solicitar reservas y registrar sus pagos. La
participación de los clientes no genera un proceso paralelo, porque sus
solicitudes ingresan al mismo flujo de validación y de registro que las
atendidas por el personal; esto permite aplicar un criterio uniforme de
disponibilidad y de confirmación. Se trata de un sistema de información
operativo orientado a la gestión comercial, no de un sistema gerencial ni de
un CRM independiente, aunque incluya funciones de registro y seguimiento de
clientes.

\section*{Campo de acción}

El campo de acción se concentra en el control de disponibilidad y la
confirmación de reservas dentro del proceso de venta de paquetes turísticos de
Vizcacha Tours, Potosí. Abarca la administración de paquetes, salidas,
clientes y pasajeros, así como el registro de cotizaciones y pagos, sin
extenderse a la contabilidad general ni a los servicios de otras empresas. La
interacción del cliente se limita a la consulta, a la solicitud de reserva y al
seguimiento de su propia reserva.

\capituloSinNumero{Objetivos}
\label{cap:objetivos}

\section*{Objetivo general}

Desarrollar un sistema web de gestión de reservas que mejore la consistencia
del proceso de venta de paquetes turísticos en Vizcacha Tours de la
ciudad de Potosí.

\section*{Objetivos específicos}

\begin{enumerate}
  \item Fundamentar teóricamente la gestión de reservas, el control de
        disponibilidad de paquetes turísticos y los criterios de diseño de
        sistemas web relacionados con la variable dependiente del proyecto.
  \item Diagnosticar el proceso actual de atención y confirmación de reservas
        en Vizcacha Tours mediante entrevistas, observación y revisión
        autorizada de documentos, para identificar las principales
        dificultades y necesidades de la agencia.
  \item Diseñar y elaborar los diferentes componentes de un sistema web que
        permita administrar paquetes, clientes, fechas y reservas, validar los
        cupos disponibles y generar información de apoyo para mejorar el
        proceso estudiado.
\end{enumerate}

\capituloSinNumero{Población y muestra}
\label{cap:poblacion}

\section*{Población}

La población de Vizcacha Tours para el tercer trimestre de 2026 está
conformada por dos grupos. El primero es el personal que participa en el
proceso de reservas: un administrador, tres agentes de atención y ventas, dos
responsables de operaciones y dos personas de apoyo administrativo y pagos, lo
que totaliza ocho trabajadores. El segundo son los clientes que consultaron o
reservaron durante ese periodo; los registros internos de la agencia
contabilizan 50 clientes. En total, la población asciende a 58 personas. Los
clientes que realizaron reservas forman parte de este segundo grupo, por lo
que no se considera un grupo separado; los proveedores no forman parte de la
población de estudio.

El sistema previsto contempla que los clientes dispongan de una cuenta propia:
los que la utilizan consultan la disponibilidad, solicitan reservas y registran
pagos dentro del sistema, mientras que el resto continúa siendo atendido por
el personal de la agencia. Esta distinción no modifica el catálogo de la
población, porque ambos grupos quedan comprendidos en las 58 personas
registradas.

Los datos provienen de los registros internos de Vizcacha Tours
correspondientes al periodo indicado.

\section*{Muestra}

Dado que la población total es de 58 personas, cantidad menor o igual a 60, se
aplica un criterio censal. Por ello, no se calcula una muestra probabilística:
se trabajará con los ocho trabajadores y con los 50 clientes registrados en el
periodo. El censo permite recoger información de todos los participantes sin
inferencias muestrales y se ajusta al tamaño reducido de la agencia y al
propósito del diagnóstico.

\begin{tablaapa}{tab:poblacion}{Población y muestra de Vizcacha Tours, tercer trimestre de 2026}
\begin{tabular}{@{}p{7.2cm}cc@{}}
\toprule
Grupo de estudio & Población & Muestra \\
\midrule
Personal de la agencia & 8 & 8 (censo) \\
Clientes que consultaron o reservaron & 50 & 50 (censo) \\
\midrule
\textbf{Total} & \textbf{58} & \textbf{58 (censo)} \\
\bottomrule
\end{tabular}
\nota{Elaboración propia con datos de los registros internos de Vizcacha Tours (2026).}
\end{tablaapa}

\chapter{Justificación social}
\label{cap:justificacion}

El proyecto busca beneficiar principalmente a los trabajadores de Vizcacha
Tours que atienden consultas y reservas. Contar con una fuente de información
sobre los paquetes, las fechas y los cupos puede facilitar su trabajo y ayudar
a que todos sigan el mismo procedimiento antes de confirmar una solicitud. De
esta forma se pretende reducir el tiempo que se dedica a revisar datos o
solucionar diferencias entre reservas.

Los clientes también pueden resultar beneficiados. Al solicitar un viaje
esperan recibir información clara sobre la disponibilidad, el precio y los
servicios incluidos. Si la agencia maneja datos actualizados y registra
correctamente cada cambio, será más sencillo dar respuestas oportunas y
evitar confusiones al momento de confirmar una reserva. El acceso directo al
sistema refuerza este beneficio: el cliente consulta la disponibilidad y el
estado de su reserva sin depender de la atención presencial, y recibe el
aviso cuando cambia la fecha o la capacidad de una salida a la que está
vinculado.

De manera indirecta, una mejor organización de los pasajeros confirmados puede
servir para coordinar el transporte, los guías y el hospedaje. La unificación
de los canales de atención beneficia también a la agencia, porque la consulta
y la solicitud del cliente ingresan al mismo criterio de disponibilidad y de
registro que aplica el personal. Por tratarse de una propuesta académica, no
se afirmará que el sistema aumentará las ventas o reducirá gastos en una
cantidad determinada hasta que se realicen pruebas y exista información que
permita demostrarlo.

\chapter{Alcances del sistema}
\label{cap:alcances}

El proyecto comprenderá el análisis del proceso de venta y reserva de
paquetes turísticos de Vizcacha Tours y el desarrollo de un sistema web que
permita administrar paquetes, programar salidas, consultar la disponibilidad
de cupos, registrar clientes y pasajeros, elaborar cotizaciones, gestionar
reservas, registrar anticipos y saldos, administrar proveedores y generar
reportes. La prioridad será verificar los cupos y registrar correctamente las
reservas; las demás funciones se incorporarán como apoyo a ese proceso
principal. El sistema será utilizado por el personal de la agencia y por los
clientes, que accederán con una cuenta propia para consultar la información
de los paquetes y solicitar sus reservas.

\section{Módulos propuestos}

\begin{itemize}
  \item \textbf{Paquetes turísticos:} registro de destinos, descripción,
        duración, itinerario, precios y servicios incluidos.
  \item \textbf{Salidas y disponibilidad:} programación de fechas, capacidad
        de pasajeros y cálculo de cupos disponibles.
  \item \textbf{Clientes y pasajeros:} registro de datos necesarios para
        atender y dar seguimiento a las solicitudes.
  \item \textbf{Cotizaciones:} preparación y consulta de propuestas según el
        paquete solicitado y sus condiciones.
  \item \textbf{Reservas:} creación, consulta, confirmación, modificación y
        cancelación con validación de cupos. La reserva puede originarse en la
        atención del personal o en la solicitud del cliente; en ambos casos se
        registra en estado pendiente y requiere la confirmación del personal.
  \item \textbf{Pagos:} registro de anticipos y saldos, con pago en efectivo y
        simulación de banca móvil o código QR, sin integración con pasarelas
        bancarias reales. El anticipo no puede ser menor al 30\% del total de
        la reserva, la totalidad también puede pagarse por la web y el pago
        registrado por el cliente queda pendiente de validación por parte del
        personal. La devolución del dinero se rige por tramos según los días
        de anticipación de la cancelación y su ejecución queda fuera del
        sistema.
  \item \textbf{Acceso del cliente:} registro de cuenta, con el correo y el
        documento obligatorios y únicos, consulta de paquetes,
        salidas y cupos disponibles, solicitud de reserva, seguimiento de sus
        propias reservas, registro de pagos y pagos simulados, con acceso
        restringido a la información que le corresponde.
  \item \textbf{Proveedores:} datos básicos de transporte, hospedaje y guías
        vinculados a los paquetes.
  \item \textbf{Reportes y usuarios:} consultas por fecha, ocupación, reservas
        y accesos según el rol de cada usuario, sea trabajador o cliente.
\end{itemize}

Para facilitar su uso, se procurará una interfaz sencilla y adaptable a
distintos tamaños de pantalla.

\section{Canales de atención y confirmación}

El sistema contempla dos canales de atención que desembocan en el mismo
proceso de reserva. El primero es la atención presencial y telefónica, en la
que el personal registra la consulta, prepara la cotización y confirma la
reserva. El segundo es el acceso directo del cliente al sistema web, en el que
consulta los paquetes, las salidas y los cupos disponibles, y remite su
solicitud de reserva junto con el registro de su pago.

En ambos casos la reserva se crea en estado pendiente y requiere la validación
y la confirmación del personal, con la misma verificación de disponibilidad y
el mismo registro de auditoría. La confirmación exige el respaldo del anticipo
validado, y las solicitudes del cliente que no son confirmadas dentro del
plazo de 24 horas quedan vencidas y liberan automáticamente los cupos
reservados. De esta manera, el acceso directo del
cliente amplía la cobertura del proceso sin duplicar reglas ni alterar las
tareas de confirmación, seguimiento y control que corresponden al personal.

\section{Avisos al cliente}

Los cambios en la fecha o en la capacidad de una salida Programada y las
confirmaciones de reserva se comunicarán al cliente afectado a través de las
pantallas del sistema, lo que reduce la demora en comunicar estas novedades y
mantiene informada a la persona que viaja.

\capituloSinNumero{Límites del proyecto}
\label{cap:limites}

El desarrollo se limitará a Vizcacha Tours y a los paquetes que administra
directamente. No incluirá conexión automática con los sistemas de aerolíneas
ni con sistemas globales de distribución. Tampoco se desarrollará una
aplicación móvil nativa, aunque el sistema podrá abrirse desde un navegador de
celular.

No se implementarán facturación tributaria ni contabilidad completa. El
sistema permitirá registrar pagos en efectivo y simular pagos por banca móvil
o código QR, con fines de seguimiento administrativo, sin conectarse a
pasarelas bancarias reales. Para el desarrollo y las pruebas se utilizarán
datos ficticios o autorizados, de modo que no se exponga información personal
de los clientes. Se aplicarán medidas básicas de control de acceso y
protección de datos, tomando como referencia las recomendaciones de OWASP
Foundation \parencite{owasp2025}.

\section*{Límites del acceso directo del cliente}

El acceso del cliente al sistema se limita a la consulta de la información de
los paquetes y de la disponibilidad, a la solicitud de reservas, al
seguimiento de sus propias reservas y al registro de pagos. El sistema no
confirma reservas por sí solo: la verificación de la disponibilidad, la
validación del pago y la confirmación de la reserva corresponden al personal de
la agencia, de manera que el pago registrado por el cliente no produce por sí
solo la confirmación de la reserva.

El cliente podrá presentar una sola solicitud activa por salida. Las
solicitudes que no son confirmadas dentro del plazo establecido quedan
vencidas y liberan de manera automática los cupos reservados. Las
modificaciones de una reserva ya confirmada, tales como el cambio de número de
pasajeros o de fecha, las realiza el personal de la agencia.

La cancelación solicitada por el cliente se registra como una solicitud y su
ejecución corresponde al personal, quien aplica los criterios de penalidad y
de devolución. La devolución de dinero por una cancelación se gestiona fuera
del sistema. El acceso directo del cliente complementa la atención del
personal y no la sustituye: las personas que no disponen de cuenta o que
prefieren la atención presencial siguen siendo atendidas en la oficina.

La apertura de una cuenta para el cliente exige datos de identificación
básicos y aplica las mismas medidas de control de acceso y de protección de
datos exigidas al personal.


%% --------------------------- Finales ---------------------------------
\iniciarSeccionesFinales
\clearpage
\printbibliography[heading=bibintoc,title={Referencias}]

\clearpage
\preliminar{Glosario}

\begin{description}[style=nextline,leftmargin=1.27cm,labelindent=0pt,%
  labelsep=0.4em,font=\bfseries,itemsep=0.25\baselineskip]
  \item[Cotización] Propuesta económica que detalla el precio, los
        servicios incluidos y las condiciones de un paquete turístico.
  \item[Cupo] Cantidad de lugares disponibles que quedan en una salida
        turística luego de descontar las reservas confirmadas.
  \item[Disponibilidad] Condición que indica si un paquete o una salida
        puede venderse en una fecha determinada.
  \item[GDS] Sistema global de distribución que permite consultar
        horarios, tarifas y disponibilidad de varios proveedores
        turísticos desde un mismo punto de acceso.
  \item[Itinerario] Secuencia de actividades, traslados y servicios
        previstos durante un viaje.
  \item[Paquete turístico] Combinación de transporte, alojamiento y
        actividades que se ofrece al cliente a un precio determinado.
  \item[Reserva] Registro de un cliente y de los pasajeros que ocuparán
        los cupos de una salida turística.
  \item[Solicitud de reserva] Registro de una reserva creado en el sistema
        por el cliente, que queda pendiente de la validación y de la
        confirmación del personal de la agencia.
  \item[Canal de atención] Vía por la que se recibe una consulta o una
        solicitud de reserva: la atención del personal de la agencia o el
        acceso directo del cliente al sistema web.
  \item[Sistema web] Aplicación a la que se accede mediante un navegador
        y que centraliza la información en un servidor.
  \item[Trazabilidad] Capacidad de seguir el historial y el estado de
        cada reserva o cambio realizado en el sistema.
  \item[Vencimiento de una solicitud] Finalización del plazo previsto para
        confirmar una reserva pendiente, que libera los cupos reservados y
        queda registrada en el historial.
\end{description}

\section*{Acrónimos}

\begin{description}[style=nextline,leftmargin=1.27cm,labelindent=0pt,%
  labelsep=0.4em,font=\bfseries,itemsep=0.15\baselineskip]
  \item[GDS] \emph{Global Distribution System}
  \item[INE] Instituto Nacional de Estadística
  \item[ISO] \emph{International Organization for Standardization}
  \item[OWASP] \emph{Open Worldwide Application Security Project}
  \item[TPS] \emph{Transaction Processing System}
  \item[UPDS] Universidad Privada Domingo Savio
\end{description}

\appendix
\clearpage

%% En los anexos el número de sección usa la letra del anexo (A.1, B.1...).
\renewcommand{\thesection}{\thechapter.\arabic{section}}

\begin{landscape}
\chapter{Árbol de problemas}
\label{anexo:arbol}
\thispagestyle{paisaje}

\setcounter{figure}{0}
\renewcommand{\thefigure}{A\arabic{figure}}

\captionsetup{hypcap=false}
\centering
\captionof{figure}{Árbol de problemas del proceso de venta de paquetes turísticos}
\label{fig:arbol}

\vspace{4pt}
\scalebox{0.85}{%% =====================================================================
%%  arbol-problemas.tex
%%  Diagrama de árbol de problemas (TikZ).
%%  Se incluye desde el Anexo A en una página horizontal.
%% =====================================================================
\begin{tikzpicture}[
  nodo/.style={
    draw=black!65, line width=0.5pt, rounded corners=1pt,
    align=center, font=\scriptsize, text width=4cm,
    inner sep=3pt, minimum height=1.7cm, fill=white
  },
  central/.style={
    nodo, text width=17cm, font=\small\bfseries,
    fill=black!8, minimum height=1.3cm
  },
  flecha/.style={
    -{Latex[length=1.8mm,width=1.4mm]},
    draw=black!65, line width=0.5pt
  },
  rotulo/.style={
    rotate=90, font=\tiny\bfseries, text=black!70, align=center
  }
]

%% ------------------------- Efectos indirectos ------------------------
\node[nodo] (I1) at (-8.4,0) {Pérdida potencial de clientes y oportunidades de venta};
\node[nodo] (I2) at (-2.8,0) {Disminución de la confianza en el servicio ofrecido};
\node[nodo] (I3) at ( 2.8,0) {Mayores costos de coordinación y corrección};
\node[nodo] (I4) at ( 8.4,0) {Dificultades en la planificación de salidas turísticas};

%% -------------------------- Efectos directos -------------------------
\node[nodo] (D1) at (-8.4,-2.6) {Demoras en la confirmación de reservas};
\node[nodo] (D2) at (-2.8,-2.6) {Errores en las reservas, cupos duplicados o mal asignados};
\node[nodo] (D3) at ( 2.8,-2.6) {Información contradictoria sobre la disponibilidad};
\node[nodo] (D4) at ( 8.4,-2.6) {Reprogramaciones, reclamos y cancelaciones};

%% --------------------------- Problema central ------------------------
\node[central] (C) at (0,-5.2) {Ineficiencia en el proceso de venta de paquetes turísticos en Vizcacha Tours, Potosí};

%% --------------------------- Causas directas -------------------------
\node[nodo] (CD1) at (-8.4,-7.8) {Verificación no uniforme de cupos antes de confirmar una reserva};
\node[nodo] (CD2) at (-2.8,-7.8) {Desfase en la actualización de ventas entre distintos canales de atención};
\node[nodo] (CD3) at ( 2.8,-7.8) {Cambios de fechas o capacidad que no se reflejan oportunamente en las reservas};
\node[nodo] (CD4) at ( 8.4,-7.8) {Criterios poco claros para modificar, cancelar y liberar reservas};

%% ------------------------- Causas indirectas -------------------------
\node[nodo] (CI1) at (-8.4,-10.4) {Responsabilidades de verificación y actualización sin definición suficiente};
\node[nodo] (CI2) at (-2.8,-10.4) {Información de salidas y reservas sin una fuente compartida y actualizada};
\node[nodo] (CI3) at ( 2.8,-10.4) {Falta de controles para identificar cambios y prevenir conflictos de capacidad};
\node[nodo] (CI4) at ( 8.4,-10.4) {Procedimientos de seguimiento y validación insuficientemente estandarizados};

%% ------------------------------ Flechas ------------------------------
%% Efecto directo -> efecto indirecto
\foreach \d/\i in {D1/I1,D2/I2,D3/I3,D4/I4} {
  \draw[flecha] (\d.north) -- (\i.south);
}

%% Problema central -> efectos directos
\foreach \d/\x in {D1/-8.4,D2/-2.8,D3/2.8,D4/8.4} {
  \draw[flecha] (C.north -| \x,0) -- (\d.south);
}

%% Causas directas -> problema central
\foreach \c/\x in {CD1/-8.4,CD2/-2.8,CD3/2.8,CD4/8.4} {
  \draw[flecha] (\c.north) -- (C.south -| \x,0);
}

%% Causas indirectas -> causas directas
\foreach \c/\d in {CI1/CD1,CI2/CD2,CI3/CD3,CI4/CD4} {
  \draw[flecha] (\c.north) -- (\d.south);
}

%% --------------------------- Rótulos de nivel -------------------------
\node[rotulo] at (-10.95,0)     {E. INDIRECTOS};
\node[rotulo] at (-10.95,-2.6)  {E. DIRECTOS};
\node[rotulo] at (-10.95,-5.2)  {PROBLEMA};
\node[rotulo] at (-10.95,-7.8)  {C. DIRECTAS};
\node[rotulo] at (-10.95,-10.4) {C. INDIRECTAS};

\end{tikzpicture}
}

\vspace{2pt}
\begin{minipage}{0.88\linewidth}
  \nota{Fuente: elaboración propia a partir del diagnóstico preliminar.}
\end{minipage}
\end{landscape}



\end{document}
